%%vsgtu709

%

% %%%

%%%   Самарский государственный технический университет

%%%   Кафедра прикладной математики и информатики

%%%

%%%   Преамбула для статьи в журнал "Вестник Самарского государственного

%%%   технического университета. Серия: «Физико-математические науки»"

%%%   Ориентирована под MikTEx 2.4 for Win32

%%%

%%%   Хотя сам журнал собирается под GNU/Linux на дистрибутиве TeXLive



\documentclass[11pt,a4paper,twoside]{article}



\usepackage[cp1251]{inputenc}

\usepackage[T2A]{fontenc}

\usepackage[english,russian]{babel}

\usepackage{samgtubib}

\usepackage[dvips]{graphicx}

\usepackage[multiple]{footmisc}



\usepackage{samgtu}

\renewcommand{\VestnikYear}{2009} % Год издания Вестника

\renewcommand{\VestnikNumber}{2\,(19)} % Номер Вестника

\renewcommand{\VestnikMonth}{Ноябрь} % Месяц выхода Вестника

\renewcommand{\VestnikData}{01.11.09} % Дата подписания Вестника в печать

\renewcommand{\VestnikUPL}{26,64} % Усл. печ. л.

\renewcommand{\VestnikUIL}{25,73} % Уч.-изд. л.

\renewcommand{\VestnikRegNumber}{479/09} % Регистрационный номер

\renewcommand{\VestnikOrder}{994} % Номер заказа







%%

% \begin{document}



\renewcommand{\firstpage}{167}

\renewcommand{\lastpage}{170}





\udk{517.95}%Дифференциальные уравнения с частными производными

\msc{35L52}%  Initial value problems for second-order hyperbolic systems



\title{Решение задачи Коши для системы уравнений Эйлера"--~Пуассона"--~Дарбу}{Решение задачи Коши для системы уравнений Эйлера"--~Пуассона"--~Дарбу}



\author{Е.\,А.~Максимова}{М\,а\,к\,с\,и\,м\,о\,в\,а~Е.\,А.}



\institute{\samgtu.}





\email{E-mail}{katyuha\_mak@mail.ru} 



\otherinf{\fullauthor{Екатерина Алексеевна Максимова} \position{аспирант} \dept{каф. прикладной математики и~информатики}}



\keywords{метод Римана, задача Коши, система уравнений Эйлера"--~Пуассона"--~Дарбу}



\workabstract{Рассмотрена система уравнений Эйлера"--~Пуассона"--~Дарбу. 

Получено решение задачи Коши для случая, когда  характеристические числа матрицы"=коэффициента комплексно"=сопряжённые  с~действительной частью из интервала $(-1/2, 0)$.}



\engtitle{Solution of the Cauchy problem for system of Euler--Poisson--Darboux equations}



\engauthor{E.\,A.~Maksimova}{M\,a\,k\,s\,i\,m\,o\,v\,a ~E.\,A.}





\enginstitute{\engsamgtu.}



\otherenginfm{\fullauthor{Ekaterina A. Maksimova}

\position{Postgraduate Student} \dept{Dept. of Applied Mathematics \& Computer Science}}



\engabstract{The system of Euler--Poisson--Darboux equations is considered, the Cauchy problem is  solved for the case, when characteristic numbers of matrix--coefficient are complex conjugate and having real part in the interval $(-1/2, 0)$.}



\engkeywords{Riemann method, the Cauchy problem, partial differential equations, Euler--Poisson--Darboux equation}



\date{21/III/2011} % Дата поступления статьи в редакцию.

\revisiondate{23/VIII/2011} % В окончательном виде



%%%%%%%%%%%%%%%%%%%%%%%%%%%%%%%%%%%%%%%%%%%%%%%%%%%%%%%%%%%%%%%%%%%%%%%%%%%%%%%%%%%%





\maketitle \small







Рассмотрим систему дифференциальных уравнений в частных производных

\begin{equation}

  \label{max:eq1}

  \frac{\partial^{2}U}{\partial x^2}-\frac{\partial^2U}{\partial y^2}-\frac{2G}{y}\frac{\partial{U}}{\partial y }=0,

\end{equation}

где 

$

  U=\left(

 u_1(x,y),

u_2(x,y)

\right)^\top 

$ "--- неизвестная вектор-функция, $G$ "--- действительная \linebreak \mbox{$2\times2$}"=матрица.



В работе  \cite{mix:dis} построена матрица Римана и с~её помощью получено решение задачи Коши для системы \eqref{max:eq1} в случае, когда спектр матрицы $G$ принадлежит интервалу  $(-1/2, 1/2)$. 

В~\cite{mix:my} получено решение задачи Коши для случая, когда собственные значения матрицы $G$ "--- комплексно"=сопряжённые числа c действительной частью из интервала $(0, 1/2)$.



Цель данной работы "--- найти решение задачи Коши для случая, когда матрица $G$ имеет комплексно"=сопряжённые собственные значения $\lambda_1$, $\lambda_2$ c действительной частью из интервала $(-1/2, 0)$:

$\lambda_1=\mu_1+i \mu_2$, 

$\lambda_2=\mu_1-i \mu_2$, 

$\mu_1 \in (-1/2, 0)$,

$\mu_2 \neq 0$.





\begin{newthm}{Задача Коши} Найти вектор-функцию $U(x, y),$ удовлетворяющую следующим условиям\/{\rm:}

\begin{enumerate}

\item[\rm 1)] $U(x, y)\in C(\bar{D})\cap C^2(D),$ где $D=\{(x,y) : 0<-y<x<y+1\};$

\item[\rm 2)] $U(x, y)$ удовлетворяет системе \rm \eqref{max:eq1}; \it

\item[\rm 3)] выполняются начальные условия

\begin{gather}

  \label{max:kr1}

  U(x, 0)=\tau(x)=(\tau_1(x),\tau_2(x))^\top,\quad  x \in [0, 1];

\\

  \label{max:kr2}

  \lim_{y \to -0}K(y)\frac{\partial U}{\partial{y}}=\nu(x)=(\nu_1(x),\nu_2(x))^\top,\quad  x \in (0, 1),

\end{gather}

где $$\displaystyle K(y)=(-y)^{2G}=(-y)^{2\mu_1}\Bigl(E\cos(2\mu_2 \ln(-y))- \frac{G-\mu_1 E}{\mu_2} \sin(2\mu_2 \ln(-y))\Bigr).$$

\end{enumerate}



\end{newthm}



В характеристических координатах 

 \begin{equation}

 \label{max:char}

  \xi=x+y,\quad   \eta=x-y

\end{equation}

область $D$ переходит в область $H=\{(\xi, \eta) : 0<\xi<\eta<1\}$, а система (\ref{max:eq1}) редуцируется к системе уравнений Эйлера"--~Пуассона"--~Дарбу специального вида:

\begin{equation}

  \label{max:eq_char}

  \frac{\partial^{2}U}{\partial \xi \partial \eta }+

  \frac{\partial U}{\partial \xi} \frac{G}{\eta-\xi}-

  \frac{\partial U}{\partial \eta} \frac{G}{\eta-\xi}=0,

\end{equation}

при этом начальные условия принимают следующий вид:

\begin{gather}

\label{max:kr1_char}

 U(\xi,\xi)=\tau(\xi),\quad \xi \in [0, 1],

\\

\label{max:kr2_char}

 \lim_{\eta \to \xi +0} K \left( \frac{\xi-\eta}{2} \right)\left(\frac{\partial U}{\partial \xi}-\frac{\partial U}{\partial \eta}\right)=\nu(\xi),\quad \xi \in (0, 1).

\end{gather}



Из свойств функции от матрицы \cite{mix:gant} следует, что функция"=матрица Римана  \linebreak

$R(\xi, \eta; \xi_0, \eta_0)$ есть вещественная матрица даже в том случае, 

когда характеристические числа матрицы $G$ комплексно"=сопряжённые.

 Запишем её аналитический вид:

\begin{multline*}

%\label {max:R}

R=V^{\mu_1} \Bigl(

\frac{G-\mu_1 E}{\mu_2}\bigl(\psi(\lambda_1,r)\cos (\mu_2 \ln V)

+\varphi(\lambda_1,r) \sin(\mu_2 \ln V) \bigr)+ 

\\

+  E\bigl(\varphi(\lambda_1,r)\cos (\mu_2 \ln V) 

+\psi(\lambda_1,r) \sin(\mu_2 \ln V)

\bigr)

\Bigr).

\end{multline*}

Здесь

\begin{gather*}

V=\frac{(\eta-\xi)^2}{(\eta-\xi_0)(\eta_0-\xi)},\quad  r=-\frac{(\xi-\xi_0)(\eta-\eta_0)}{(\xi-\eta_0)(\xi_0-\eta)},

\notag

\\

\varphi(\lambda_1,r)= \text{Re}\Bigl({}_2F_1\Bigl( 

\begin{array}{c}

 \mu_1+i \mu_2, \mu_1+i\mu_2 \\

1

\end{array};

r

\Bigr) \Bigr),  

\\

  \psi(\lambda_1,r)= \text{Im} \Bigl({}_2F_1\Bigl( 

\begin{array}{c}

 \mu_1+i \mu_2, \mu_1+i\mu_2 \\

1

\end{array};

r

\Bigr) \Bigr).

\end{gather*}



Если $U(\xi,\eta)$ является решением системы уравнений (\ref{max:eq_char}), а $R(\xi,\eta;\xi_0,\eta_0)$ "--- матрица Римана этой системы уравнений, то, используя свойства матрицы Римана и векторный аналог тождества Грина \cite{mix:bic}, получаем

\begin{multline}

\label{max:U_eps}

 U_{\varepsilon}(\xi_0,\eta_0)=\frac{1}{2}(RU)\Bigr|_{\scriptsize{

\begin{array}{l}

\xi=\eta_0-\varepsilon\\

\eta=\eta_0 

 \end{array} }}+

\frac{1}{2}(RU)\Bigr|_{\scriptsize{

\begin{array}{l}

\xi=\xi_0\\

\eta=\xi_0 +\varepsilon

 \end{array} }} +

 \\

+\frac{1}{2}\int_{\xi_0}^{\eta_0-\varepsilon}R\left(\frac{\partial U}{\partial \eta}-\frac{\partial U}{\partial \xi}\right)\Bigr|_{\eta=\xi+\varepsilon} d\xi - \\ - \frac{1}{2}\int_{\xi_0}^{\eta_0-\varepsilon} \left(

\frac{\partial R}{\partial \eta} - \frac{\partial R}{\partial \xi} + \frac{4RG}{\xi-\eta}

\right)U\Bigr|_{\eta=\xi+\varepsilon}d\xi=\sum\limits_{k=1}^{4}J_{k}(\varepsilon).

\end{multline}



В равенстве (\ref{max:U_eps}) непосредственно переходить к пределу при $\varepsilon\rightarrow 0$ нельзя, так как все внеинтегральные члены стремятся к~бесконечности, а интеграл $J_4(\varepsilon)$ не существует. 

Интегрируя по частям  $J_4(\varepsilon)$, применяя формулу автотрансформации и переходя к пределу при $\varepsilon\to 0$, получим решение задачи Коши (\ref{max:kr1_char}), (\ref{max:kr2_char}) для системы уравнений (\ref{max:eq_char}) в области $H$, которое имеет вид

\begin{multline}

\label{max:result}

U(\xi,\eta)=-2^{2 \mu_1 -1}

\left(

K_1 \psi(\lambda_1,1)+E \varphi(\lambda_1,1)

\right)

\int_{\xi}^{\eta}  \sigma^{-\mu_1}(t) \cos 

  \left(

    \mu_2 \ln \frac{4}{\sigma(t)}

  \right)\nu(t)dt-

\\

-2^{2 \mu_1 -1}

\left(

K_1 \varphi(\lambda_1,1)-E \psi(\lambda_1,1)

\right)

\int_{\xi}^{\eta} \sigma^{-\mu_1}(t) \sin 

  \left(

    \mu_2 \ln \frac{4}{\sigma(t)}

  \right)

\nu(t)dt +

\\

+(\eta-\xi)^{-2\mu_1-1} 

 \bigg[

\frac{1}{2} 

\left( 

 E\psi(-\lambda_1,1) - K_1 \varphi(-\lambda_1,1)

\right) \times

\\

\times

\int_{\xi}^{\eta}

\sigma^{\mu_1}(t)

\sin \left(

    \mu_2 \ln \frac{(\xi-\eta)^2}{\sigma(t)}

  \right)(\eta+\xi-2t)

\tau'(t)dt -

\\

- \frac{1}{2} 

\left( 

 E \varphi(-\lambda_1,1)+ K_1\psi(-\lambda_1,1)

\right)

\int_{\xi}^{\eta}

\sigma^{\mu_1}(t)

\cos \left(

    \mu_2 \ln \frac{(\xi-\eta)^2}{\sigma(t)}

  \right)

(\eta+\xi-2t)

\tau'(t)dt+

\\

+ \left( 

K_2\varphi(-\lambda_1,1)+

K_3\psi(-\lambda_1,1)

\right)

\int_{\xi}^{\eta}

\sigma^{\mu_1}(t)

\cos \left(

    \mu_2 \ln \frac{(\xi-\eta)^2}{\sigma(t)}

  \right)

\tau(t)dt+

\\

+\left(

 K_2\psi(-\lambda_1,1)-

K_3\varphi(-\lambda_1,1)

\right)

\int_{\xi}^{\eta}

\sigma^{\mu_1}(t)

\sin \left(

    \mu_2 \ln \frac{(\xi-\eta)^2}{\sigma(t)}

  \right)

\tau(t)dt

\bigg],

\end{multline}

где

\begin{gather*}

K_1=\frac{G-\mu_1 E}{\mu_2},\quad K_2=2G+E, \\

K_3=(1+2\mu_1)K_1-2\mu_2 E, \quad  \sigma(t)=(t-\xi)(\eta-t), \\

\varphi(\lambda_1,1)=1+\sum\limits_{n=1}^{\infty} \dfrac{\varphi_n(\lambda_1,1)}{(1)_n n!}, \quad  \psi(\lambda_1,1)=\sum\limits_{n=1}^{\infty} \dfrac{\psi_n(\lambda_1,1)}{(1)_n n!},

\\

\varphi_{n+1}(\lambda_1,1)=((\mu_1+n)^2-\mu_2^2)\varphi_{n}(\lambda_1,1)-2\mu_2(\mu_1+n)\psi_n(\lambda_1,1), 

\\

\psi_{n+1}(\lambda_1,1)=((\mu_1+n)^2-\mu_2^2)\psi_{n}(\lambda_1,1)+2\mu_2(\mu_1+n)\varphi_n(\lambda_1,1),

\\

\varphi_1(\lambda_1,1)=\mu_1^2-\mu_2^2,\quad  \psi_1(\lambda_1,1)=2 \mu_1 \mu_2.

\end{gather*}



Используя (\ref{max:char}), можно записать решение $U(x,y)$ в области $D$. Таким образом, справедлива следующая  теорема.



\begin{theorem} Если $\tau(x)\in C^{3} [0, 1]$ и $\nu(x)\in C^{2} (0, 1),$ то задача Коши \rm \eqref{max:kr1}, \eqref{max:kr2} \it  для уравнения \eqref{max:eq1} корректна по Адамару\rm.

\end{theorem}



\begin{newthm}{Замечание} 

Положив в \eqref{max:result} $\mu_1=0,$ получим решение задачи Коши для случая мнимых собственных значений матрицы $G.$

\end{newthm}



\vspace{-3mm}



\begin{thebibliography}{9}



\RBibitem{mix:dis}

\by Андреев~А.\,А.

\paper Об одном классе систем дифференциальных уравнений гиперболического типа 

\inbook Дифференциальные уравнения в частных производных

\bookinfo Cб. тр. мат. кафедр пединститутов РСФСР. Вып.~16

\yr 1980 

\publaddr Рязань

\publ Рязан. гос. пед. инст.

\pages 9–14

%MR0632553 

\misctransl 

\by Andreev~A.\,A. 

\paper On a class of systems of differential equations of hyperbolic type

\inbook Partial differential equations

\pages 9–14

\publ Ryazan. Gos. Ped. Inst. 

\publaddr Ryazan

\yr 1980 





\RBibitem{mix:my}

\paper Решение задачи Коши для одной системы гиперболического типа с сингулярными характеристиками

\by Андреев~А.\,А., Максимова~Е.\,А.

\publaddr Самара

\publ СамГТУ

\inbook Труды восьмой Всероссийской научной конференции с международным участием. \rm Часть~3 

\bookinfo Дифференциальные уравнения и краевые задачи

\serial Математическое моделирование и краевые задачи

\yr 2011

\pages 11--17

\misctransl

\paper The solution of the Cauchy problem for one hyperbolic system with singular characteristics

\by Andreev~A.\,A, Maksimova~E.\,A.

\inbook Proceedings of the Eighth All-Russian Scientific Conference with international participation. \rm  Part~3

\serial Matem. Mod. Kraev. Zadachi

\publaddr Samara

\publ SamGTU

\yr 2011

\pages 11--17







\RBibitem{mix:gant}

\by Гантмахер~Ф.\,Р.

\book Теория матриц

\publaddr М.

\publ Наука

\yr 1988

\totalpages 549

\misctransl

\by Gantmakher~F.\,R.

\book Theory of matrices

\publaddr Moscow

\publ Nauka

\totalpages 549







\RBibitem{mix:bic}

\by Бицадзе~А.\,В.

\book Уравнения смешанного типа

\publaddr М.

\publ Наука

\yr 1966

\totalpages 164

\transl англ. пер.:{} {} 

\by Bitsadze~A.\,V.

\book Equations of the Mixed Type 

\publ Pergamon Press

\publaddr New York

\yr 1964

\totalpages 160







\end{thebibliography}



\noindent

\hskip6.5cm \thedate \par\noindent

\hskip6.5cm\therevdate





\newpage



\chead[\fancyplain{}{\scriptsize \sl M\,a\,k\,s\,i\,m\,o\,v\,a~E.\,A.}]{\fancyplain{}{\scriptsize \sl  Solution of the Cauchy problem for system of Euler--Poisson--Darboux equations}}



\makeengtitlem



\newpage



%\makecontents

% Тестовая комманда для создания оглавления в трудах конференции

%\end{document}

